\documentclass[manuscript,screen,review]{acmart}

\AtBeginDocument{%
  \providecommand\BibTeX{{%
    Bib\TeX}}}

\setcopyright{acmlicensed}
\copyrightyear{2026}
\acmYear{2026}
\acmDOI{XXXXXXX.XXXXXXX}
\acmConference[AH'26]{Augmented Humans 2026}{March 16--19,
  2026}{Okinawa, Japan}
\acmISBN{978-1-4503-XXXX-X/2018/06}


\usepackage{framed}
\usepackage{changepage}


\begin{document}

\title{Beyond Cognitive Offloading: Conversational AI for Personal Growth}


\settopmatter{printacmref=false}

\author{Dániel Szabó}
\orcid{0009-0003-7299-9116}
\affiliation{%
  \institution{University of Oulu}
  \city{Oulu}
  \country{Finland}
}
\email{daniel.szabo@oulu.fi}

\author{Niels van Berkel}
\orcid{0000-0001-5106-7692}
\affiliation{%
  \institution{Aalborg University}
  \city{Aalborg}
  \country{Denmark}
}
\email{nielsvanberkel@cs.aau.dk}

\author{Benjamin Tag}
\orcid{0000-0002-7831-2632}
\affiliation{%
  \institution{University of New South Wales}
  \city{Sydney}
  \country{Australia}
}
\email{benjamin.tag@unsw.edu.au}

\author{Rune Møberg Jacobsen}
\orcid{0000-0002-1877-1845}
\affiliation{%
  \institution{Aalborg University}
  \city{Aalborg}
  \country{Denmark}
}
\email{runemj@cs.aau.dk}

\author{Jorge Goncalves}
\orcid{0000-0002-0117-0322}
\affiliation{%
  \institution{University of Melbourne}
  \city{Melbourne}
  \country{Australia}
}
\email{jorge.goncalves@unimelb.edu.au}

\author{Xinrui Fang}
\orcid{0000-0001-5444-8722}
\affiliation{%
  \institution{The University of Tokyo}
  \city{Tokyo}
  \country{Japan}
}
\email{xinrui.fang@iis-lab.org}

\author{Zefan Sramek}
\orcid{0000-0003-4541-0399}
\affiliation{%
  \institution{The University of Tokyo}
  \city{Tokyo}
  \country{Japan}
}
\email{zefans@iis-lab.org}

\author{Koji Yatani}
\orcid{0000-0003-4192-0420}
\affiliation{%
  \institution{The University of Tokyo}
  \city{Tokyo}
  \country{Japan}
}
\email{koji@iis-lab.org}


\author{Aku Visuri}
\orcid{0000-0001-7127-4031}
\affiliation{%
  \institution{University of Oulu}
  \city{Oulu}
  \country{Finland}
}
\email{aku.visuri@oulu.fi}


\author{Simo Hosio}
\orcid{0000-0002-9609-0965}
\affiliation{%
  \institution{University of Oulu}
  \city{Oulu}
  \country{Finland}
}
\email{simo.hosio@oulu.fi}



\renewcommand{\shortauthors}{Szabó et al.}

\begin{abstract}
Conversational AI now helps people perform a plethora of everyday tasks, especially online. 
Such cognitive offloading can enhance task performance, but does not always contribute to the development of the underlying human abilities. 
As a result, there are increasing concerns about the potential deleterious effects of overusing AI in one's work. 
This workshop explores a timely and positive framing: cognitive augmentation, where conversational AI can help people grow rather than just perform their tasks.
We bring together a diverse set of researchers to work on how these systems might foster inherently human cognitive capabilities, such as meta-cognition, self-motivation, and emotional intelligence, without creating dependency.
\end{abstract}

%%
%% The code below is generated by the tool at http://dl.acm.org/ccs.cfm.
%% Please copy and paste the code instead of the example below.
%%
\begin{CCSXML}
<ccs2012>
   <concept>
       <concept_id>10003120.10003121</concept_id>
       <concept_desc>Human-centered computing~Human computer interaction (HCI)</concept_desc>
       <concept_significance>500</concept_significance>
       </concept>
 </ccs2012>
\end{CCSXML}

\ccsdesc[500]{Human-centered computing~Human computer interaction (HCI)}
%%
%% Keywords. The author(s) should pick words that accurately describe
%% the work being presented. Separate the keywords with commas.
\keywords{Cognitive offloading, cognitive augmentation, conversational AI}


%%
%% This command processes the author and affiliation and title
%% information and builds the first part of the formatted document.
\maketitle

\section{Introduction}

%Digital technology has enabled us to enhance many aspects of everyday life. Virtually unlimited storage capacity allows us to augment memory and revisit recorded events. Complex arithmetic that would take hours to solve by hand can be completed at the press of a button on a calculator. Seamless automated translation now enables conversations across language barriers with just earphones. 
Conversational AI (CAI) is now available widely online. 
This has led many to use online tools such as Claude and ChatGPT for various everyday tasks that earlier required manual work.
This is a form of \textit{cognitive offloading}~\cite{Risko2016}, i.e. delegating work that would normally require thinking, freeing up our cognitive resources for other purposes. 
%This, however, also bears the risk of reducing the need for direct cognitive engagement. 
%Increased efficiency and reduced errors are clear benefits of these developments. 
Yet, by diminishing sustained engagement with the underlying cognitive tasks, such offloading may limit opportunities for development of related skills, and in some cases, even contribute to their decline ~\cite{yourbrainongpt25}. 
%Yet, while they may augment our ability to perform tasks, they do not support our personal skills development. 

%The calculator completes the arithmetic, but this does not necessarily improve our mathematical competence. %we do not become better at mathematics. 
%A translation feature bridges the language gap, but it does not help us gain proficiency in the language. Similarly, in the domain of well-being, digital interventions often focus on short-term symptom alleviation or eliciting positive emotions, rather than training users to achieve those states independently~\cite{Mitsea2024}. 
%With the increased focus of CAI on inherently human aspects of life, be it to support our thinking and decision-making or to help us deal with complex emotions, 
This workshop probes how this novel technology can cultivate our abilities rather than, over time, deteriorate them. There is increasing concern over AI potentially harming our cognitive abilities~\cite{yourbrainongpt25}, and consequently, researchers are increasingly looking into mitigating this~\cite{yatani2024aiextrahericsfosteringhigherorder}.
%In doing so, this workshop seeks to explore how conversational AI can help us develop our minds, our cognitive faculties, rather than perform tasks for us. 
We adopt the term `\textit{cognitive augmentation}' to distinguish our goal from cognitive offloading: Where offloading delegates tasks to technology, augmentation strengthens the person, while retaining the benefits of efficiency and reduced errors.
%CAI can adapt to individual users and engage in reflective dialogue, both of which are critical for learning and growth. Recent advances in large language models (LLM) have made such interactions increasingly natural and accessible, and 
Many CAI users now also engage in use that is not just functional, and deals e.g. with learning and personal growth~(\textit{e.g,}~\cite{Wester2025SelfCare, Skjuve2021}). Yet we lack a clear understanding of how to design these systems to achieve such personal development. 
%Specifically, we aim to investigate how these systems can be designed to foster personal growth—supporting users in developing meta-cognitive skills, emotional intelligence, and internal motivation—without creating dependency.
The workshop addresses several interrelated questions. What constitutes cognitive augmentation in the context of CAI? How might CAI improve the users' \textit{meta-cognitive skills}? Can AI interactions help people gain \textit{emotional intelligence}? What role might such systems play in cultivating \textit{internal motivation}? Can AI support users' self-development of cognitive capabilities beyond AI interactions? And critically, what are the risks of getting this wrong?

% The workshop addresses several interrelated questions:
% \begin{itemize}
%     \item What constitutes cognitive augmentation in the context of conversational AI?
%     \item How might conversational AI improve the users' \textit{meta-cognitive skills}?
%     \item Can AI interactions help people gain \textit{emotional intelligence}?
%     \item What role might such systems play in cultivating \textit{internal motivation}?
%     \item Can AI support the self-development of users' conscious control beyond the AI interactions?
%     \item And critically, what are the risks of getting this wrong?
% \end{itemize}


% discussion/reading -> AI interaction -> discussion/reading?

% - building resilience / slowing down the decline
% - improving / recovering

% - cognitive tool

% Inoculating helps to avoid decline
% skeptical, less open, 

\section{Background}

%\subsection{The Challenge of AI Over-reliance}

CAI is the current most widely adopted AI tool; for example, the ChatGPT alone handles over 2 billion daily queries and attracts over 800 million monthly active visitors worldwide~\cite{demandsage_chatgpt_stats}. 
While such tools offer unprecedented efficiency and accessibility, their widespread adoption also poses risks to personal growth: over-reliance on AI, where users increasingly delegate cognitive tasks to external AI tools, leads to cognitive offloading~\cite{Risko2016, sparrow2011google}, that may hinder the development of self-regulatory skills~\cite{yourbrainongpt25, Gerlich_2025}.
% In particular, excessive cognitive offloading and over-reliance on digital assistants can hinder personal growth and the development of self-regulatory skills~\cite{yourbrainongpt25, Gerlich_2025}. 
For example, Gerlich et al.~\cite{Gerlich_2025} demonstrated that over-reliance on AI may undermine critical thinking skills: specifically, younger participants in their study exhibited higher reliance on AI and lower critical thinking scores.
Lee et al. \cite{lee2025the} found that higher confidence in generative AI systems correlates with reduced critical thinking, whereas higher self-confidence in one’s own knowledge promotes it.
Meanwhile, many existing digital mental health tools also reflect this tension. 
Such tools primarily focus on short-term symptom management or immediate relief, rather than scaffolding long-term psychological resilience, growth or cognitive capability~\cite{guo2024large, zohar2024an}. 
For example, prior work has raised concerns that over-reliance on AI in mental health contexts may hinder therapists’ and patients in-depth engagement between each others~\cite{ zohar2024an}.
These limitations motivates a shift from AI tools that replace human cognition to tools that actively support and strengthen it. 

%\subsection{Cognitive Tools for Augmentation}
Rather than automating cognitive tasks, prior research on \emph{Cognitive Tools} (or Cognitive Aids) conceptualises digital systems as generalizable supports that engage users in cognitive processing rather than performing it on their behalf~\cite{cognitive_tools_for_learning_1992}. Such tools function as both mental and computational devices, structuring attention, prompting reflection, and extending users' thinking processes~\cite{flexible_cognitive_tools_1990}. They can deliver targeted ``Cognitive Boosts'' that improve decision-making~\cite{hertwig2017nudging} or provide ``Cognitive Inoculation''—a form of psychological immunization against external influences such as misinformation~\cite{shin2024algorithmic}.
With the proliferation of LLMs, recent research has demonstrated the use of CAI as \emph{Cognitive Tools} to promote personal growth, including enhancing higher-order thinking skills~\cite{tanprasert2024debate, wambsganss2021argueTutor}, mitigating confirmation bias~\cite{zhang2024see}, and supporting mental well-being~\cite{geng2025beyond, wester2025using}.
For example, Tanprasert et al.~\cite{tanprasert2024debate} designed a debate chatbot that helps users improve critical thinking while watching YouTube videos and Geng et al.~\cite{geng2025beyond} developed a multi-chatbot system and found that it improved women’s coping motivation for premenstrual syndrome (PMS) management.
These ideas are particularly relevant in the context of psychological wellbeing, which is increasingly understood as a multi-faceted construct encompassing relationships, emotions, autonomy and personal growth~\cite{Dr_M_Dhanabhakyam_2023}.

%\subsection{Conversational AI and Cognitive Augmentation}
% These ideas are particularly relevant in the context of psychological wellbeing, which is increasingly understood as a multi-faceted construct encompassing relationships, emotions, autonomy and personal growth~\cite{Dr_M_Dhanabhakyam_2023}. 
At the same time, recent technological developments, especially the integration of AI tools in social and interactive media platforms, have been associated with reduced face-to-face interaction and altered social dynamics~\cite{Nortei2025}. As we integrate AI more deeply into our lives, understanding its immediate as well as longer-term impact on these fundamental aspects of wellbeing becomes critical.
CAI, when thoughtfully designed as cognitive tools for positive psychology interventions, can provide personalised guidance for practice of cognitive re-framing, goal-setting and strengths identification. This can encourage users to view challenges as growth opportunities, fostering positive behavioural changes~\cite{narain2020promoting}.

However, without careful design constraints, these same systems risk encouraging over-dependence, potentially undermining motivation, engagement, and effort in training meta-skills required for self-consciousness, self-regulation, and adaptability \cite{richards2023principlist}. 
This tension highlights a central design challenge for CAI: how can it support users in cultivating positive mental states and cognitive skills without becoming a substitute for those capacities? Addressing this question requires reframing CAI not as an external source of regulation or insight, but as a scaffold that gradually enables individuals to self-induce, sustain, and generalise these cognitive and emotional processes over time.
Indeed, while disconnective practices have received some attention in the context of mobile communication technologies~\cite{mannell2018typology}, considering how and when users might best disconnect from CAI so as to rely on their internalised abilities remains an open challenge.
% positive mental states, presenting a design challenge for CAs aiming to augment one's cognitive growth.

%\vspace{-8px}

\subsection{Target Constructs}
%This workshop specifically focuses on how CAI can be designed to foster specific human capabilities rather than replacing them. 
%Drawing on recent research on supporting skills related to our own conscious control of our development (e.g. ~\cite{Mitsea2024}), 
We identify three primary constructs for exploration based on recent work (see e.g.~\cite{Mitsea2024}). 
\textbf{(1) Meta-cognition:} The ability to be self-conscious, self-regulated, and adaptive in one's thinking. Using GenAI tools imposes significant strain on meta-cognition \cite{10.1145/3613904.3642902}, and in this workshop, we explore how CAI can support the training of these meta-skills, encouraging users to reflect on their cognitive processes rather than offloading them entirely;
\textbf{(2) Self-motivation:} The internal drive to engage with tasks and sustain effort. We examine how CAI can instead motivate users to identify personal strengths, set realistic goals, and perceive challenges as growth opportunities. This also includes examining how users frame learning as goals, as opposed to simply offloading tasks~\cite{Li14042025};
\textbf{(3) Emotional Intelligence:} The capacity to understand and manage emotions. Recent concerns highlight how CAI can create pseudo-intimacy that might hamper emotional capabilities~\cite{10.3389/fpsyg.2024.1410462}. We investigate how CAs can go beyond facilitating positive mental states by scaffolding users' emotional intelligence. 

%Together, these constructs represent the core pillars of cognitive augmentation that this workshop aims to explore.

\section{Content of Workshop}

This workshop will be in-person only to encourage networking and the low-friction sharing of ideas, prioritising critical discourse and collaboration. 
%The activities are structured to move from \textit{theoretical framing} to \textit{practical application} and finally to \textit{critical reflection}.

% \subsection{Opening Provocation}
% We will kick off with a keynote presentation titled "???????????????". This talk challenges the prevailing design ethos of "frictionless" interaction, arguing that by smoothing away all cognitive effort, we risk atrophying the very muscles required for conscious control. The keynote will frame the day's challenge: how to design AI that supports \textit{Meta-cognition}, fosters \textit{Self-motivation}, and scaffolds \textit{Emotional Intelligence}, all while focusing on fostering lasting self-development efforts over directly eliciting the positive behaviours. It posits that current approaches often prioritize convenience, when we actually need "gym equipment" to strengthen the mind.

\paragraph*{Lightning Pitches into Grand Challenges}
Participants submit their own ideas of interest and items that are relevant to the workshop topic, which will be briefly presented in a rapid session. These ideas are organised into themes during the session, after which we will conduct an activity to formulate a set of grand challenges for cognitive augmentation using conversational AI. Using simple affinity diagramming, we cluster the key challenges. Then, we facilitate a rapid prototyping session to brainstorm initial solutions and designs for addressing the challenges.
The prototypes can be "vibe coded" using solutions like Claude or Google Antigravity. 
The organisers are also experienced with paper prototyping and will provide participants with printed templates that can be used without digital tools, so everyone can participate even without digital development experience.
%Groups will be able to choose between various delivery methods, e.g., "vibe coding" or Wizard of Oz simulations, to prototype their concepts.



%\begin{itemize}
%    \item \textbf{Round 1:} Groups are formed around subjects related to the accepted paper submissions and the workshop's topics, and they spend 30 minutes discussing specific challenges. This will be followed by a 15-minute session where groups present their key discussion points to the other groups orally. 
%\item \textbf{Round 2:} Building on the first round, groups will spend another 30 minutes discussing potential design solutions or frameworks. This round also concludes with a 15-minute presentation session.
%\end{itemize}

%\paragraph*{Practical Prototyping}
%Following the structured discussions, participants will engage in a 30-minute hands-on activity. 
%\paragraph*{Critique and Synthesis}
%We will circle back to the three core problems identified earlier, Meta-cognition, Self-motivation, and Emotional Intelligence, to evaluate whether the proposed designs successfully scaffold individuals' internal capabilities for Conscious Control, or merely provide a temporary solution.

\paragraph*{Logistics and Requirements}
We require a standard workshop room with movable furniture to facilitate group work and an easy-to-use standard A/V system for presentations. Stable internet access is critical as participants will be using cloud-based LLMs. 

\section{Goals and Outcomes}

The primary goal of this workshop is to collaboratively explore and address key questions and challenges related to the core concept of cognitive augmentation through conversational AI with practical solutions. Specific outcomes include:

\noindent\textbf{Prototype Repository:} A web-based gallery on the workshop home page of the "Vibe Coded" prototypes and design concepts produced during the workshop, serving as tangible examples of how to use CAI for cognitive augmentation.
 
\noindent\textbf{Joint Position Paper:} A collaborative publication that summarises the workshop's findings, outlines a future research agenda, and presents potential results of any future studies based on the workshop's prototypes.

\noindent\textbf{Community Building:} Establishing a network of researchers interested in the intersection of AI, psychology, and education, specifically focused on the target constructs of meta-cognition, self-motivation, and emotional intelligence.

We will set up a website for the workshop to host all submissions and outputs. The organisers will advertise the workshop, including on HCI-related mailing lists.

\section{Schedule}

The workshop is planned for 3.5 hours, after which the organisers will coordinate a voluntary social event.


% Option A, 3-stage plan - 3 hours 30 minutes 
% 1. Discussions about design issues
% 2. Discussions about approaches of solutions
% 3. Practical demonstrations of potential solutions
%\noindent\small\begin{tabular}{ll}
%    Introduction and Welcome & (15 min) \\
%    Self-introductions and Lightning Pitches of Accepted Works & (45 min)  \\
%    Coffee Break & (15 min) \\
%    Plenary Discussion & (30 min) \\
%    Practical Prototyping in Small Groups & (60 min) \\
%    Synthesis, Final Discussion and Wrap-up & (30 min) \end{tabular}

\noindent\small\begin{tabular}{ll}
    Welcome and Introduction to the Topic & (10 min) \\
    Self-introductions and Lightning Pitches of Accepted Submissions & (20 min)  \\
    Panel Discussion & (30min) \\
    Collection and Formulation of Grand Challenge Questions & (30 min) \\
    Coffee Break & (15 min) \\
    Practical Prototyping in Small Groups & (75 min) \\
    Synthesis, Final Discussion and Wrap-up & (30 min) \end{tabular}


% Option B, submission-focused plan - 3 hours 30 minutes 
% 1. All submissions are very shortly presented and then discussed in depth
% 2. At the end, we form groups. Either one group per accepted paper or one group per 3 conscious control subjects (meta-cognition, emotional intelligence, self-motivation)
% \noindent\small\begin{tabular}{ll} 
%     Introduction and Welcome & (10 min) \\
%     Self-introductions & (5 min)  \\
%     Keynote / Opening Provocation & (15 min) \\
%     Accepted Work 1 presentation & (5 min) \\
%     Accepted Work 1 discussion & (15 min) \\
%     Accepted Work 2 presentation & (5 min) \\
%     Accepted Work 2 discussion & (15 min) \\
%     Coffee Break & (15 min) \\
%     Accepted Work 3 presentation & (5 min) \\
%     Accepted Work 3 discussion & (15 min) \\
%     Accepted Work 4 presentation & (5 min) \\
%     Accepted Work 4 discussion & (15 min) \\
%     Forming Groups and Group Discussions & (30 min) \\
%     Synthesis, Final Discussion and Wrap-up & (30 min) \end{tabular}
    


%Please describe how many attendees the workshop will have, who you will invite to submit papers to your workshop, and how you will reach/recruit them, including a draft call for participation. How will you conduct the process of reviewing and accepting submissions? Please also describe how you will support an accessible and inclusive workshop. Due to room capacity, workshops will be limited to a maximum of 20 on-site participants, including organisers.

\section{Recruitment and Reviewing}
%participants
The workshop will host up to 20 in-person participants (including organisers). We invite submissions from researchers, practitioners, and students in HCI, Human-AI Interaction, Psychology, Education and related fields, as described in the draft of the participation call below. Recruitment will occur through direct invitation and public call. 
%We will coordinate with accepted authors so that all workshop papers are published on arXiv.

%review
% Papers will undergo double-blind review by organisers, evaluated on relevance, originality, clarity and potential for discussion.


\begin{framed}
\footnotesize

% We invite those interested in the workshop topic to submit papers (2-4 pages) outlining their perspectives, research, or design concepts related to cognitive augmentation. Topics of interest include, but are not strictly limited to:
% \begin{itemize}
%     \item{Design of conversational AI for meta-cognition,}
%     \item{AI-supported interventions for emotional intelligence,}
%     \item{Technologies fostering self-motivation,}
%     \item{Ethical considerations in cognitive augmentation.}
% \end{itemize}
% Submissions will be reviewed based on relevance, originality, clarity, and potential for sparking discussion. Accepted submissions will be presented during the lightning pitches and discussed in the structured group discussions. Participants will also have the opportunity to engage in practical prototyping sessions, contributing to a web-based gallery of prototypes and a joint position paper. 

We invite submissions that identify a core problem or propose a solution related to the risks of cognitive offloading in conversational AI, including (but not limited to): \textit{Dependency and reduced self-reflection}, \textit{Erosion of critical thinking or meta-cognition}, \textit{Over-automation of emotional or moral reasoning}, \textit{Design strategies that foster learning, agency, or long-term growth}.

Participants may submit work in any format of their choice, including: \textit{Paper or short essay}, \textit{Poster}, \textit{Video (accessible online}, \textit{Interactive website, app, prototype, or tool (submit a brief PDF including the URL and a short description)}, \textit{Speculative or critical design artifact}. 
Our goal here is to be maximally open to all formats, and will be happy to accommodate author requests in this matter.
We welcome research, design, artistic, and practice-based contributions. 
Submissions should clearly articulate the problem addressed and/or the proposed approach, and authors must be ready to pitch their submission in the short lightning talks session at the beginning of the workshop.

The workshop aims to create a multidisciplinary space for discussion, critique, and hands-on exploration of how conversational AI can move beyond cognitive offloading toward meaningful personal growth.
\end{framed}

% accessibility and inclusivity
To ensure an inclusive and accessible workshop, we will ask participants to share any related requests, and do our best to accommodate all in cooperation with the conference organisers. 
Further, we will use inclusive language in all communications and materials. The workshop organisers are highly experienced in successful workshops, having organised several in key HCI venues such as CHI, CSCSW, and UbiComp, ensuring the workshop atmosphere will be safe and encouraging for participants.

\section*{Workshop organisers}

\noindent\textbf{Dániel Szabó} is currently pursuing a PhD at the University of Oulu, Finland. His primary research interest is conversational human-AI interaction.

\noindent\textbf{Niels van Berkel} is a professor at Aalborg University, Denmark. His current projects focus on mental health and decision support.

\noindent\textbf{Benjamin Tag} is a senior lecturer at the University of New South Wales in Sydney, Australia. His research focuses on affective computing, human factors, and ubiquitous computing. 

\noindent\textbf{Rune Møberg Jacobsen} is a postdoc at Aalborg University, Denmark. His research focuses on AI-supported Decision-Making and Conversational Interfaces for health. 

\noindent\textbf{Xinrui Fang} is a PhD student at the University of Tokyo, Japan. His research focuses on exploring how AI can be designed to foster human critical thinking skills and support mental well-being.

\noindent\textbf{Zefan Sramek} is a PhD student at the University of Tokyo, Japan. His current research is on the relationship between social media use, agency, and psychological well-being.

\noindent\textbf{Jorge Goncalves} is an associate professor a the University of Melbourne. Dr. Goncalves has organised workshops across CHI, CSCW and UbiComp conference and has a strong background in applied LLM and HCI research.

\noindent\textbf{Koji Yatani} is a professor at the University of Tokyo. Dr. Yatani has a strong background in ubiquitous computing and, more recently, on conversational AI systems especially in the context of human wellbeing. Dr. Yatani has organised several workshops in various HCI conferences and leads the Interactive Intelligent Systems Laboratory.

\noindent\textbf{Aku Visuri} is an Academy Fellow at the University of Oulu and has a strong background in HCI and recommendation systems.

\noindent\textbf{Simo Hosio} is a professor at the University of Oulu. Dr. Hosio has organised workshops at CHI, CSCW and UbiComp, and he leads the Crowd Computing research group that focuses on crowdsourcing and applied AI systems for human wellbeing.

%%
%% The next two lines define the bibliography style to be used, and
%% the bibliography file.
\bibliographystyle{ACM-Reference-Format}
\bibliography{REFERENCES}

%%
%% If your work has an appendix, this is the place to put it.
\appendix


\end{document}
\endinput
%%
%% End of file `sample-manuscript.tex'.
